
\documentclass{article}
\usepackage{colortbl}
\usepackage{makecell}
\usepackage{multirow}
\usepackage{supertabular}

\begin{document}

\newcounter{utterance}

\centering \large Interaction Transcript for game `matchit\_ascii', experiment `similar\_grid\_1', episode 3 with claude{-}opus{-}4{-}8{-}azure.
\vspace{24pt}

{ \footnotesize  \setcounter{utterance}{1}
\setlength{\tabcolsep}{0pt}
\begin{supertabular}{c@{$\;$}|p{.15\linewidth}@{}p{.15\linewidth}p{.15\linewidth}p{.15\linewidth}p{.15\linewidth}p{.15\linewidth}}
    \# & $\;$A & \multicolumn{4}{c}{Game Master} & $\;\:$B\\
    \hline

    \theutterance \stepcounter{utterance}  
    & & \multicolumn{4}{p{0.6\linewidth}}{
        \cellcolor[rgb]{0.9,0.9,0.9}{
            \makecell[{{p{\linewidth}}}]{
                \texttt{\tiny{[P1$\langle$GM]}}
                \texttt{Tá tú ag glacadh páirte i gcluiche tomhais comhoibríoch.} \\
\texttt{Is é an sprioc a fháil amach an bhfuil an greille seo agus pictiúr eile nach féidir liomsa amháin a fheiceáil, mar an gcéanna.} \\
\texttt{Déan cur síos ar do ghreille ar dtús, le do thoil. Ansin, tabharfaidh mé mo chur síos agus is féidir linn ceisteanna a chur ar a chéile faoi na greillí chun a dhéanamh amach an bhfuil siad mar an gcéanna. Anois cuir tús le do chur síos gairid ar an ngreille le "TUAIRISC:" agus an cur síos ina dhiaidh. Ná cuir aon rud eile leis. Seo é do ghreille:} \\
\texttt{X X X X X} \\
\texttt{X ▢ ▢ ▢ ▢} \\
\texttt{X X X ▢ ▢} \\
\texttt{X ▢ ▢ ▢ ▢} \\
\texttt{X ▢ ▢ ▢ ▢} \\
            }
        }
    }
    & \\ \\

    \theutterance \stepcounter{utterance}  
    & \multicolumn{4}{p{0.6\linewidth}}{
        \cellcolor[rgb]{0.9,0.9,0.9}{
            \makecell[{{p{\linewidth}}}]{
                \texttt{\tiny{[P1$\rangle$GM]}}
                \texttt{TUAIRISC: Is greille 5x5 í seo. Tá an chéad ró líonta go hiomlán le X (cúig X). Sa dara ró, tá X sa chéad chill amháin agus tá na ceithre chill eile folamh. Sa tríú ró, tá X sna chéad trí chill agus tá an dá chill dheireanacha folamh. Sa cheathrú ró, tá X sa chéad chill amháin agus tá na ceithre chill eile folamh. Sa chúigiú ró, tá X sa chéad chill amháin agus tá na ceithre chill eile folamh. Mar sin, tá an chéad cholún ar fad líonta le X, tá an chéad ró ar fad líonta le X, agus tá dhá X breise sa tríú ró (an dara agus an tríú cill).} \\
            }
        }
    }
    & & \\ \\

    \theutterance \stepcounter{utterance}  
    & & & \multicolumn{2}{p{0.3\linewidth}}{
        \cellcolor[rgb]{0.9,0.9,0.9}{
            \makecell[{{p{\linewidth}}}]{
                \texttt{\tiny{[GM$|$GM]}}
                \texttt{continue} \\
            }
        }
    }
    & & \\ \\

    \theutterance \stepcounter{utterance}  
    & & \multicolumn{4}{p{0.6\linewidth}}{
        \cellcolor[rgb]{0.9,0.9,0.9}{
            \makecell[{{p{\linewidth}}}]{
                \texttt{\tiny{[GM$\rangle$P2]}}
                \texttt{Tá tú ag glacadh páirte i gcluiche tomhais comhoibríoch.} \\
\texttt{Is é an sprioc a fháil amach an bhfuil an greille seo agus pictiúr eile nach féidir liomsa amháin a fheiceáil, mar an gcéanna.} \\
\texttt{Déan cur síos ar do ghreille ar dtús, le do thoil. Ansin, tabharfaidh mé mo chur síos agus is féidir linn ceisteanna a chur ar a chéile faoi na greillí chun a dhéanamh amach an bhfuil siad mar an gcéanna. Anois cuir tús le do chur síos gairid ar an ngreille le "TUAIRISC:" agus an cur síos ina dhiaidh. Ná cuir aon rud eile leis. Seo é do ghreille:} \\
\texttt{X ▢ ▢ ▢ ▢} \\
\texttt{X ▢ ▢ ▢ ▢} \\
\texttt{X X X ▢ ▢} \\
\texttt{X ▢ ▢ ▢ ▢} \\
\texttt{X X X X X} \\
            }
        }
    }
    & \\ \\

    \theutterance \stepcounter{utterance}  
    & & & \multicolumn{4}{p{0.6\linewidth}}{
        \cellcolor[rgb]{0.9,0.9,0.9}{
            \makecell[{{p{\linewidth}}}]{
                \texttt{\tiny{[GM$\langle$P2]}}
                \texttt{TUAIRISC: Greille 5x5. Tá an chéad cholún ar chlé líonta go hiomlán (na cúig chill). Tá an cúigiú sraith (an tsraith ag bun) líonta go hiomlán freisin. Sa tríú sraith, tá na chéad trí chill líonta. Tá gach cill eile folamh. Tá cuma na litreach "E" ar an bpatrún.} \\
            }
        }
    }
     \\ \\

    \theutterance \stepcounter{utterance}  
    & & & \multicolumn{2}{p{0.3\linewidth}}{
        \cellcolor[rgb]{0.9,0.9,0.9}{
            \makecell[{{p{\linewidth}}}]{
                \texttt{\tiny{[GM$|$GM]}}
                \texttt{continue} \\
            }
        }
    }
    & & \\ \\

    \theutterance \stepcounter{utterance}  
    & & \multicolumn{4}{p{0.6\linewidth}}{
        \cellcolor[rgb]{0.9,0.9,0.9}{
            \makecell[{{p{\linewidth}}}]{
                \texttt{\tiny{[GM$\rangle$P2]}}
                \texttt{Seo é mo CHUR SÍOSTUAIRISC: Is greille 5x5 í seo. Tá an chéad ró líonta go hiomlán le X (cúig X). Sa dara ró, tá X sa chéad chill amháin agus tá na ceithre chill eile folamh. Sa tríú ró, tá X sna chéad trí chill agus tá an dá chill dheireanacha folamh. Sa cheathrú ró, tá X sa chéad chill amháin agus tá na ceithre chill eile folamh. Sa chúigiú ró, tá X sa chéad chill amháin agus tá na ceithre chill eile folamh. Mar sin, tá an chéad cholún ar fad líonta le X, tá an chéad ró ar fad líonta le X, agus tá dhá X breise sa tríú ró (an dara agus an tríú cill).} \\
\texttt{Anois cuir ceist chun gnéithe nua de mo ghreille a fháil amach a d'fhéadfadh a bheith difriúil ó do ghreille. Tosaigh le "CEIST:" agus ná cuir aon rud eile leis.} \\
            }
        }
    }
    & \\ \\

    \theutterance \stepcounter{utterance}  
    & & & \multicolumn{4}{p{0.6\linewidth}}{
        \cellcolor[rgb]{0.9,0.9,0.9}{
            \makecell[{{p{\linewidth}}}]{
                \texttt{\tiny{[GM$\langle$P2]}}
                \texttt{CEIST: An bhfuil an chéad ró (an tsraith ag barr) líonta go hiomlán le cúig X, nó an bhfuil an chéad chill amháin líonta inti?} \\
            }
        }
    }
     \\ \\

    \theutterance \stepcounter{utterance}  
    & & & \multicolumn{2}{p{0.3\linewidth}}{
        \cellcolor[rgb]{0.9,0.9,0.9}{
            \makecell[{{p{\linewidth}}}]{
                \texttt{\tiny{[GM$|$GM]}}
                \texttt{continue} \\
            }
        }
    }
    & & \\ \\

    \theutterance \stepcounter{utterance}  
    & & \multicolumn{4}{p{0.6\linewidth}}{
        \cellcolor[rgb]{0.9,0.9,0.9}{
            \makecell[{{p{\linewidth}}}]{
                \texttt{\tiny{[P1$\langle$GM]}}
                \texttt{Seo é mo CHUR SÍOSTUAIRISC: Greille 5x5. Tá an chéad cholún ar chlé líonta go hiomlán (na cúig chill). Tá an cúigiú sraith (an tsraith ag bun) líonta go hiomlán freisin. Sa tríú sraith, tá na chéad trí chill líonta. Tá gach cill eile folamh. Tá cuma na litreach "E" ar an bpatrún.} \\
\texttt{CEIST: An bhfuil an chéad ró (an tsraith ag barr) líonta go hiomlán le cúig X, nó an bhfuil an chéad chill amháin líonta inti?Tosaigh do fhreagra le "FREAGRA:" agus ná cuir aon rud eile leis.} \\
            }
        }
    }
    & \\ \\

    \theutterance \stepcounter{utterance}  
    & \multicolumn{4}{p{0.6\linewidth}}{
        \cellcolor[rgb]{0.9,0.9,0.9}{
            \makecell[{{p{\linewidth}}}]{
                \texttt{\tiny{[P1$\rangle$GM]}}
                \texttt{FREAGRA: Tá an chéad ró (an tsraith ag barr) líonta go hiomlán le cúig X i mo ghreille.} \\
            }
        }
    }
    & & \\ \\

    \theutterance \stepcounter{utterance}  
    & & & \multicolumn{2}{p{0.3\linewidth}}{
        \cellcolor[rgb]{0.9,0.9,0.9}{
            \makecell[{{p{\linewidth}}}]{
                \texttt{\tiny{[GM$|$GM]}}
                \texttt{continue} \\
            }
        }
    }
    & & \\ \\

    \theutterance \stepcounter{utterance}  
    & & \multicolumn{4}{p{0.6\linewidth}}{
        \cellcolor[rgb]{0.9,0.9,0.9}{
            \makecell[{{p{\linewidth}}}]{
                \texttt{\tiny{[P1$\langle$GM]}}
                \texttt{Anois cuir ceist chun gnéithe nua de mo ghreille a fháil amach a d'fhéadfadh a bheith difriúil ó do ghreille. Tosaigh le "CEIST:" agus ná cuir aon rud eile leis.} \\
            }
        }
    }
    & \\ \\

    \theutterance \stepcounter{utterance}  
    & \multicolumn{4}{p{0.6\linewidth}}{
        \cellcolor[rgb]{0.9,0.9,0.9}{
            \makecell[{{p{\linewidth}}}]{
                \texttt{\tiny{[P1$\rangle$GM]}}
                \texttt{CEIST: I do ghreille, an bhfuil an cúigiú sraith (an tsraith ag bun) líonta go hiomlán le cúig chill, nó an bhfuil an chéad chill amháin líonta inti?} \\
            }
        }
    }
    & & \\ \\

    \theutterance \stepcounter{utterance}  
    & & & \multicolumn{2}{p{0.3\linewidth}}{
        \cellcolor[rgb]{0.9,0.9,0.9}{
            \makecell[{{p{\linewidth}}}]{
                \texttt{\tiny{[GM$|$GM]}}
                \texttt{continue} \\
            }
        }
    }
    & & \\ \\

    \theutterance \stepcounter{utterance}  
    & & \multicolumn{4}{p{0.6\linewidth}}{
        \cellcolor[rgb]{0.9,0.9,0.9}{
            \makecell[{{p{\linewidth}}}]{
                \texttt{\tiny{[GM$\rangle$P2]}}
                \texttt{FREAGRA: Tá an chéad ró (an tsraith ag barr) líonta go hiomlán le cúig X i mo ghreille.} \\
\texttt{CEIST: I do ghreille, an bhfuil an cúigiú sraith (an tsraith ag bun) líonta go hiomlán le cúig chill, nó an bhfuil an chéad chill amháin líonta inti?Tosaigh do fhreagra le "FREAGRA:" agus ná cuir aon rud eile leis.} \\
            }
        }
    }
    & \\ \\

    \theutterance \stepcounter{utterance}  
    & & & \multicolumn{4}{p{0.6\linewidth}}{
        \cellcolor[rgb]{0.9,0.9,0.9}{
            \makecell[{{p{\linewidth}}}]{
                \texttt{\tiny{[GM$\langle$P2]}}
                \texttt{FREAGRA: I mo ghreille, tá an cúigiú sraith (an tsraith ag bun) líonta go hiomlán le cúig X.} \\
            }
        }
    }
     \\ \\

    \theutterance \stepcounter{utterance}  
    & & & \multicolumn{2}{p{0.3\linewidth}}{
        \cellcolor[rgb]{0.9,0.9,0.9}{
            \makecell[{{p{\linewidth}}}]{
                \texttt{\tiny{[GM$|$GM]}}
                \texttt{continue} \\
            }
        }
    }
    & & \\ \\

    \theutterance \stepcounter{utterance}  
    & & \multicolumn{4}{p{0.6\linewidth}}{
        \cellcolor[rgb]{0.9,0.9,0.9}{
            \makecell[{{p{\linewidth}}}]{
                \texttt{\tiny{[GM$\rangle$P2]}}
                \texttt{Anois cuir ceist chun gnéithe nua de mo ghreille a fháil amach a d'fhéadfadh a bheith difriúil ó do ghreille. Tosaigh le "CEIST:" agus ná cuir aon rud eile leis.} \\
            }
        }
    }
    & \\ \\

    \theutterance \stepcounter{utterance}  
    & & & \multicolumn{4}{p{0.6\linewidth}}{
        \cellcolor[rgb]{0.9,0.9,0.9}{
            \makecell[{{p{\linewidth}}}]{
                \texttt{\tiny{[GM$\langle$P2]}}
                \texttt{CEIST: I do ghreille, an bhfuil an chéad ró (an tsraith ag barr) líonta go hiomlán le cúig X, nó an bhfuil an chéad chill amháin líonta inti?} \\
            }
        }
    }
     \\ \\

    \theutterance \stepcounter{utterance}  
    & & & \multicolumn{2}{p{0.3\linewidth}}{
        \cellcolor[rgb]{0.9,0.9,0.9}{
            \makecell[{{p{\linewidth}}}]{
                \texttt{\tiny{[GM$|$GM]}}
                \texttt{continue} \\
            }
        }
    }
    & & \\ \\

    \theutterance \stepcounter{utterance}  
    & & \multicolumn{4}{p{0.6\linewidth}}{
        \cellcolor[rgb]{0.9,0.9,0.9}{
            \makecell[{{p{\linewidth}}}]{
                \texttt{\tiny{[P1$\langle$GM]}}
                \texttt{FREAGRA: I mo ghreille, tá an cúigiú sraith (an tsraith ag bun) líonta go hiomlán le cúig X.} \\
\texttt{CEIST: I do ghreille, an bhfuil an chéad ró (an tsraith ag barr) líonta go hiomlán le cúig X, nó an bhfuil an chéad chill amháin líonta inti?Tosaigh do fhreagra le "FREAGRA:" agus ná cuir aon rud eile leis.} \\
            }
        }
    }
    & \\ \\

    \theutterance \stepcounter{utterance}  
    & \multicolumn{4}{p{0.6\linewidth}}{
        \cellcolor[rgb]{0.9,0.9,0.9}{
            \makecell[{{p{\linewidth}}}]{
                \texttt{\tiny{[P1$\rangle$GM]}}
                \texttt{FREAGRA: I mo ghreille, tá an chéad ró (an tsraith ag barr) líonta go hiomlán le cúig X.} \\
            }
        }
    }
    & & \\ \\

    \theutterance \stepcounter{utterance}  
    & & & \multicolumn{2}{p{0.3\linewidth}}{
        \cellcolor[rgb]{0.9,0.9,0.9}{
            \makecell[{{p{\linewidth}}}]{
                \texttt{\tiny{[GM$|$GM]}}
                \texttt{continue} \\
            }
        }
    }
    & & \\ \\

    \theutterance \stepcounter{utterance}  
    & & \multicolumn{4}{p{0.6\linewidth}}{
        \cellcolor[rgb]{0.9,0.9,0.9}{
            \makecell[{{p{\linewidth}}}]{
                \texttt{\tiny{[P1$\langle$GM]}}
                \texttt{Anois cuir ceist chun gnéithe nua de mo ghreille a fháil amach a d'fhéadfadh a bheith difriúil ó do ghreille. Tosaigh le "CEIST:" agus ná cuir aon rud eile leis.} \\
            }
        }
    }
    & \\ \\

    \theutterance \stepcounter{utterance}  
    & \multicolumn{4}{p{0.6\linewidth}}{
        \cellcolor[rgb]{0.9,0.9,0.9}{
            \makecell[{{p{\linewidth}}}]{
                \texttt{\tiny{[P1$\rangle$GM]}}
                \texttt{CEIST: I do ghreille, sa tríú sraith, an bhfuil na chéad trí chill líonta go díreach (agus an dá chill dheireanacha folamh), nó an bhfuil líon difriúil cealla líonta inti?} \\
            }
        }
    }
    & & \\ \\

    \theutterance \stepcounter{utterance}  
    & & & \multicolumn{2}{p{0.3\linewidth}}{
        \cellcolor[rgb]{0.9,0.9,0.9}{
            \makecell[{{p{\linewidth}}}]{
                \texttt{\tiny{[GM$|$GM]}}
                \texttt{continue} \\
            }
        }
    }
    & & \\ \\

    \theutterance \stepcounter{utterance}  
    & & \multicolumn{4}{p{0.6\linewidth}}{
        \cellcolor[rgb]{0.9,0.9,0.9}{
            \makecell[{{p{\linewidth}}}]{
                \texttt{\tiny{[GM$\rangle$P2]}}
                \texttt{FREAGRA: I mo ghreille, tá an chéad ró (an tsraith ag barr) líonta go hiomlán le cúig X.} \\
\texttt{CEIST: I do ghreille, sa tríú sraith, an bhfuil na chéad trí chill líonta go díreach (agus an dá chill dheireanacha folamh), nó an bhfuil líon difriúil cealla líonta inti?Tosaigh do fhreagra le "FREAGRA:" agus ná cuir aon rud eile leis.} \\
            }
        }
    }
    & \\ \\

    \theutterance \stepcounter{utterance}  
    & & & \multicolumn{4}{p{0.6\linewidth}}{
        \cellcolor[rgb]{0.9,0.9,0.9}{
            \makecell[{{p{\linewidth}}}]{
                \texttt{\tiny{[GM$\langle$P2]}}
                \texttt{FREAGRA: I mo ghreille, sa tríú sraith, tá na chéad trí chill líonta go díreach agus tá an dá chill dheireanacha folamh.} \\
            }
        }
    }
     \\ \\

    \theutterance \stepcounter{utterance}  
    & & & \multicolumn{2}{p{0.3\linewidth}}{
        \cellcolor[rgb]{0.9,0.9,0.9}{
            \makecell[{{p{\linewidth}}}]{
                \texttt{\tiny{[GM$|$GM]}}
                \texttt{continue} \\
            }
        }
    }
    & & \\ \\

    \theutterance \stepcounter{utterance}  
    & & \multicolumn{4}{p{0.6\linewidth}}{
        \cellcolor[rgb]{0.9,0.9,0.9}{
            \makecell[{{p{\linewidth}}}]{
                \texttt{\tiny{[GM$\rangle$P2]}}
                \texttt{Anois cuir ceist chun gnéithe nua de mo ghreille a fháil amach a d'fhéadfadh a bheith difriúil ó do ghreille. Tosaigh le "CEIST:" agus ná cuir aon rud eile leis.} \\
            }
        }
    }
    & \\ \\

    \theutterance \stepcounter{utterance}  
    & & & \multicolumn{4}{p{0.6\linewidth}}{
        \cellcolor[rgb]{0.9,0.9,0.9}{
            \makecell[{{p{\linewidth}}}]{
                \texttt{\tiny{[GM$\langle$P2]}}
                \texttt{CEIST: I do ghreille, sa dara sraith, an bhfuil an chéad chill amháin líonta le X agus na ceithre chill eile folamh?} \\
            }
        }
    }
     \\ \\

    \theutterance \stepcounter{utterance}  
    & & & \multicolumn{2}{p{0.3\linewidth}}{
        \cellcolor[rgb]{0.9,0.9,0.9}{
            \makecell[{{p{\linewidth}}}]{
                \texttt{\tiny{[GM$|$GM]}}
                \texttt{continue} \\
            }
        }
    }
    & & \\ \\

    \theutterance \stepcounter{utterance}  
    & & \multicolumn{4}{p{0.6\linewidth}}{
        \cellcolor[rgb]{0.9,0.9,0.9}{
            \makecell[{{p{\linewidth}}}]{
                \texttt{\tiny{[P1$\langle$GM]}}
                \texttt{FREAGRA: I mo ghreille, sa tríú sraith, tá na chéad trí chill líonta go díreach agus tá an dá chill dheireanacha folamh.} \\
\texttt{CEIST: I do ghreille, sa dara sraith, an bhfuil an chéad chill amháin líonta le X agus na ceithre chill eile folamh?Tosaigh do fhreagra le "FREAGRA:" agus ná cuir aon rud eile leis.} \\
            }
        }
    }
    & \\ \\

    \theutterance \stepcounter{utterance}  
    & \multicolumn{4}{p{0.6\linewidth}}{
        \cellcolor[rgb]{0.9,0.9,0.9}{
            \makecell[{{p{\linewidth}}}]{
                \texttt{\tiny{[P1$\rangle$GM]}}
                \texttt{FREAGRA: I mo ghreille, sa dara sraith, tá an chéad chill amháin líonta le X agus tá na ceithre chill eile folamh.} \\
            }
        }
    }
    & & \\ \\

    \theutterance \stepcounter{utterance}  
    & & & \multicolumn{2}{p{0.3\linewidth}}{
        \cellcolor[rgb]{0.9,0.9,0.9}{
            \makecell[{{p{\linewidth}}}]{
                \texttt{\tiny{[GM$|$GM]}}
                \texttt{continue} \\
            }
        }
    }
    & & \\ \\

    \theutterance \stepcounter{utterance}  
    & & \multicolumn{4}{p{0.6\linewidth}}{
        \cellcolor[rgb]{0.9,0.9,0.9}{
            \makecell[{{p{\linewidth}}}]{
                \texttt{\tiny{[P1$\langle$GM]}}
                \texttt{Anois cuir ceist chun gnéithe nua de mo ghreille a fháil amach a d'fhéadfadh a bheith difriúil ó do ghreille. Tosaigh le "CEIST:" agus ná cuir aon rud eile leis.} \\
            }
        }
    }
    & \\ \\

    \theutterance \stepcounter{utterance}  
    & \multicolumn{4}{p{0.6\linewidth}}{
        \cellcolor[rgb]{0.9,0.9,0.9}{
            \makecell[{{p{\linewidth}}}]{
                \texttt{\tiny{[P1$\rangle$GM]}}
                \texttt{CEIST: I do ghreille, sa cheathrú sraith, an bhfuil an chéad chill amháin líonta le X agus na ceithre chill eile folamh?} \\
            }
        }
    }
    & & \\ \\

    \theutterance \stepcounter{utterance}  
    & & & \multicolumn{2}{p{0.3\linewidth}}{
        \cellcolor[rgb]{0.9,0.9,0.9}{
            \makecell[{{p{\linewidth}}}]{
                \texttt{\tiny{[GM$|$GM]}}
                \texttt{continue} \\
            }
        }
    }
    & & \\ \\

    \theutterance \stepcounter{utterance}  
    & & \multicolumn{4}{p{0.6\linewidth}}{
        \cellcolor[rgb]{0.9,0.9,0.9}{
            \makecell[{{p{\linewidth}}}]{
                \texttt{\tiny{[GM$\rangle$P2]}}
                \texttt{FREAGRA: I mo ghreille, sa dara sraith, tá an chéad chill amháin líonta le X agus tá na ceithre chill eile folamh.} \\
\texttt{CEIST: I do ghreille, sa cheathrú sraith, an bhfuil an chéad chill amháin líonta le X agus na ceithre chill eile folamh?Tosaigh do fhreagra le "FREAGRA:" agus ná cuir aon rud eile leis.} \\
            }
        }
    }
    & \\ \\

    \theutterance \stepcounter{utterance}  
    & & & \multicolumn{4}{p{0.6\linewidth}}{
        \cellcolor[rgb]{0.9,0.9,0.9}{
            \makecell[{{p{\linewidth}}}]{
                \texttt{\tiny{[GM$\langle$P2]}}
                \texttt{FREAGRA: I mo ghreille, sa cheathrú sraith, tá an chéad chill amháin líonta le X agus tá na ceithre chill eile folamh.} \\
            }
        }
    }
     \\ \\

    \theutterance \stepcounter{utterance}  
    & & & \multicolumn{2}{p{0.3\linewidth}}{
        \cellcolor[rgb]{0.9,0.9,0.9}{
            \makecell[{{p{\linewidth}}}]{
                \texttt{\tiny{[GM$|$GM]}}
                \texttt{continue} \\
            }
        }
    }
    & & \\ \\

    \theutterance \stepcounter{utterance}  
    & & \multicolumn{4}{p{0.6\linewidth}}{
        \cellcolor[rgb]{0.9,0.9,0.9}{
            \makecell[{{p{\linewidth}}}]{
                \texttt{\tiny{[GM$\rangle$P2]}}
                \texttt{Anois déan cinneadh. Cad a cheapann tú: an bhfuil do ghreille agus an ghreille eile a thuairiscítear mar an gcéanna? Scríobh "CINNEADH: greille céanna" má cheapann tú gurb í an ghreille chéanna í nó "CINNEADH: greillí éagsúla" má cheapann tú gur greillí éagsúla iad. Ná cuir aon rud eile leis.} \\
            }
        }
    }
    & \\ \\

    \theutterance \stepcounter{utterance}  
    & & & \multicolumn{4}{p{0.6\linewidth}}{
        \cellcolor[rgb]{0.9,0.9,0.9}{
            \makecell[{{p{\linewidth}}}]{
                \texttt{\tiny{[GM$\langle$P2]}}
                \texttt{CINNEADH: greillí éagsúla} \\
            }
        }
    }
     \\ \\

    \theutterance \stepcounter{utterance}  
    & & & \multicolumn{2}{p{0.3\linewidth}}{
        \cellcolor[rgb]{0.9,0.9,0.9}{
            \makecell[{{p{\linewidth}}}]{
                \texttt{\tiny{[GM$|$GM]}}
                \texttt{continue} \\
            }
        }
    }
    & & \\ \\

    \theutterance \stepcounter{utterance}  
    & & & \multicolumn{2}{p{0.3\linewidth}}{
        \cellcolor[rgb]{0.9,0.9,0.9}{
            \makecell[{{p{\linewidth}}}]{
                \texttt{\tiny{[GM$|$GM]}}
                \texttt{success} \\
            }
        }
    }
    & & \\ \\

    \theutterance \stepcounter{utterance}  
    & & \multicolumn{4}{p{0.6\linewidth}}{
        \cellcolor[rgb]{0.9,0.9,0.9}{
            \makecell[{{p{\linewidth}}}]{
                \texttt{\tiny{[P1$\langle$GM]}}
                \texttt{FREAGRA: I mo ghreille, sa cheathrú sraith, tá an chéad chill amháin líonta le X agus tá na ceithre chill eile folamh.} \\
\texttt{Anois déan cinneadh. Cad a cheapann tú: an bhfuil do ghreille agus an ghreille eile a thuairiscítear mar an gcéanna? Scríobh "CINNEADH: greille céanna" má cheapann tú gurb í an ghreille chéanna í nó "CINNEADH: greillí éagsúla" má cheapann tú gur greillí éagsúla iad. Ná cuir aon rud eile leis.} \\
            }
        }
    }
    & \\ \\

    \theutterance \stepcounter{utterance}  
    & \multicolumn{4}{p{0.6\linewidth}}{
        \cellcolor[rgb]{0.9,0.9,0.9}{
            \makecell[{{p{\linewidth}}}]{
                \texttt{\tiny{[P1$\rangle$GM]}}
                \texttt{CINNEADH: greillí éagsúla} \\
            }
        }
    }
    & & \\ \\

    \theutterance \stepcounter{utterance}  
    & & & \multicolumn{2}{p{0.3\linewidth}}{
        \cellcolor[rgb]{0.9,0.9,0.9}{
            \makecell[{{p{\linewidth}}}]{
                \texttt{\tiny{[GM$|$GM]}}
                \texttt{continue} \\
            }
        }
    }
    & & \\ \\

    \theutterance \stepcounter{utterance}  
    & & & \multicolumn{2}{p{0.3\linewidth}}{
        \cellcolor[rgb]{0.9,0.9,0.9}{
            \makecell[{{p{\linewidth}}}]{
                \texttt{\tiny{[GM$|$GM]}}
                \texttt{success} \\
            }
        }
    }
    & & \\ \\

\end{supertabular}
}

\end{document}
